\clearpage
\section*{Second owner-verdict response: one primary question}
\subsection*{What does the original benchmark measure?}
The paper's primary question is whether the labels and split support an antimicrobial-function claim. They do not: review status is coupled to class in the original cohort. An AUROC can be numerically reproducible and still answer the wrong biological question. The 0.921 test number summarizes ranking on the constructed UniProt collection; 0.926 summarizes APD6 positives versus reviewed UniProt negatives. The latter changes the positive source but retains a class-specific negative source and cannot establish clean cross-database transfer. The selected 239 reviewed positives in the original test are a small post-hoc subgroup, not a replacement training cohort. The central contribution is the audit trail that makes this failure visible and makes a valid rebuild testable.

\subsection*{Overlap with the first owner verdict and genuinely new work}
The first verdict's P0--P5 paper response already placed the audit first, treated those AUROCs as confounded, limited the GNN and candidate claims, and specified matched negatives and alignment-based family splits as future requirements. This second verdict overlaps those items; it does not represent an independent confirmation or a second clean benchmark. Its new paper-side emphasis is to keep one primary question, retract unsupported priority language such as ``first automated census,'' avoid ad hoc admissibility thresholds masquerading as validated biology, and count only directly used tools and independent study-level sources. Whether a result is novel in the literature needs a bounded search; absence of a known prior source in this project is not proof of firstness.

\subsection*{The completed secondary 0.45 sensitivity arm}
The committed \texttt{results/study21\_hard\_split\_t0.45.json} on science main (commit \texttt{621ce7e}) records ten folds and 3,821 clusters. Mean MCC is 0.8248 (fold SD 0.0296; one-sided lower 95\% 0.8077); mean AUROC is 0.9667. The fixed 0.799 hard point-value gate passes, including the reported lower confidence bound. The stronger 0.8328 mean-point gate fails. This is a secondary threshold arm, not a selection opportunity after seeing 0.40 and 0.35. The primary 0.40 result still passes both mean-point gates by a narrow margin, while 0.35 passes only the 0.799 mean-point hard gate and its lower bound is 0.7972. All three are within-dataset tests. Their agreement on a weaker gate does not remove source-review confounding or prove family-level independence.

\subsection*{Limits on flags, tools and source counts}
An operational admissibility rule is a filter chosen for a pipeline, not a calibrated probability of AMP activity. Without blinded expert review, an assay and an independently curated comparator, a flagged accession is a hypothesis. The screen was score-conditioned after training and cannot be retrospectively preregistered. Counts of external APIs, software dependencies, database records, accession files, source queries and independent biological studies are different denominators; infrastructure calls and MGNify/PubMed records should not be promoted into distinct research tools or independent cohorts. An auditable inventory must name each directly used tool and each dataset's study-level parent, retain exclusions, and mark genuinely pending external measurements as pending. No new numerical tool or cohort count is asserted here.

\subsection*{Decision path and next endpoint}
For a prospective benchmark, begin by auditing the positive and negative accession provenance before selecting a model. If review status or collection source predicts class by construction, redesign the cohort before fitting; if a negative is merely unannotated, describe it as such. Freeze reviewed-positive and organism/length-matched reviewed-negative definitions, then lock sequence-alignment family partitions and a truly external source. Report paired uncertainty at an appropriate family unit, calibration and low-FPR errors by biologically annotated class. A clean rebuild is the next scientific endpoint, not an improvement claim. [PENDING: committed reviewed-only retrain outcome, alignment-based partition and external cohort.] This revision prioritizes substantive manuscript text and does not include a rendered presentation.
